\begin{workedexample}[Worked Example: SNR scaling with field]
MRI signal-to-noise rises steeply with static field. In the
body-noise-dominated regime SNR grows roughly linearly with $B_0$, so moving from
$1.5\ \text{T}$ to $3\ \text{T}$ roughly doubles SNR. That gain drives high-field
imaging --- at the cost of higher RF frequency ($128$ vs.\ $64\ \text{MHz}$),
stronger susceptibility artefacts, and higher SAR (which scales as $B_0^2$).
\end{workedexample}

\begin{keytakeaways}
\begin{itemize}[leftmargin=1.4em,itemsep=2pt]
\item The magnet sets $B_0$ (and frequency); gradients encode space; the RF chain excites and receives.
\item Higher $B_0$ improves SNR but raises frequency, SAR ($\propto B_0^2$), and susceptibility artefacts.
\item Shimming homogenizes $B_0$; RF and gradient shielding contain stray fields.
\end{itemize}
\end{keytakeaways}

\begin{exercises}
\item Roughly how does SNR change from $1.5\ \text{T}$ to $3\ \text{T}$?
\item How does whole-body SAR scale with $B_0$?
\item Which improves with higher field: SNR or susceptibility artefacts?
\end{exercises}

\furtherreading{Nishimura~\cite{nishimura}; Haacke~\cite{haacke} (ch.~27).}
